\subsection{Complete source acquisition and local refinement}

\paragraph{A common matter and compact-field carrier.}
The full 192-CAR native source keeps all 114 source terms and 24 even pair
intertwiners. Two parallel $U(1)$ links and an integer quantum reference
solve the Gauss constraints exactly as
$E_1=3Q_1-n$, $E_2=n$, $R=-3(Q_0+Q_1)$, $n\in\mathbb Z$.
The entire physical carrier is $\mathcal F_{192}\otimes\ell^2(\mathbb Z)$.
For the inherited source $H_0$, define the declared positive extension
\[
 H=H_0+\frac{(3Q_1-n)^2+n^2}{32}
 +\frac1{64}\sum_{j=1}^{96}(d_{0j}-d_{1j})^*(d_{0j}-d_{1j})
 +\frac1{64}\left(I-\frac{S+S^*}{2}\right),\quad S|n\rangle=|n+1\rangle.
\]
It is self-adjoint on $D(n^2)$, positive, and has compact resolvent on
this finite spatial graph. All native cross-cell terms carry their exact
integer dressing. Charge, local current, reference and complete-instrument
work identities hold on every retained matter/flux sector. The constant
pair maps have zero connection; any field-dependent replacement retains
$E+\mathcal A$ and every term of its square.

The full one-lepton reducing operators have second eigenvalue bounded
from below by $1/16$ after subtracting their inherited rest coefficient.
A rational trial with exact full integer-tail residual gives
$\mu-v/(1/16-\mu)\le E_0\le\mu$ by Temple's inequality.
The inherited coloured and charged occupation floors exceed the cycle
vacuum bound $1/64$, while neutral-only states share its field operator.
This proves that the cycle vacuum is the full-source ground.
Subtracting this vacuum yields a common charged gap increase
approximately $7.535487\times10^{-6}$ GeV with the supplied electron unit.
Neutral one-lepton shifts cancel exactly. This is a calculated architecture
dependence, not a physical rest-mass prediction. For every native
preparation at $n=0$, the added mean is
$9\alpha\langle Q_1^2\rangle+t\langle K\rangle+\beta$; the operator
difference is unbounded on the full flux space. A native-process compiler
and an independent colour gauge field are not supplied by this $U(1)$
extension (\path{model/openwave_matter_flux_source}).

A separate optional normalization uses the exact computed spectral number
$\delta=E_{\rm charged}-E_v$, with $0<\delta<1/64$, and defines
$H_\delta=H-\delta N_{\rm ch}$. The inherited inequality
$H_0\ge cN_{\rm colour}+N_{\rm ch}+H_{\rm neutral}$, $c>1$, gives
$H_\delta\ge cN_{\rm colour}+(1-\delta)N_{\rm ch}+H_{\rm neutral}$.
Every coloured or charged sector remains above the cycle vacuum, and
the neutral-only sector retains exactly $H_v$ plus positive terms.
Thus the global vacuum is unchanged and the same coefficient restores
all inherited charged one-lepton rest coefficients; neutral ones are
unchanged. This is an optional source convention with its own generator,
not a native-process derivation or a new physical mass prediction.
Weak transitions can change $N_{\rm ch}$, so current and work formulas
must be evaluated for $H_\delta$.

\paragraph{Repeated finite mass-clock records.}
A finite phase reference supports two native mass-clock cycles when source reset is
implemented as an energy-conserving exchange with an explicit empty battery.
For preparation coordinates $c^2+s^2=1$, the two isolated source coherences
are $is e^{imt_1}/2$ and $isc e^{imt_2}/2$; depletion reduces contrast without
changing the ideal frequency. Complete source and resource occupations give
exact gap-switching work $m(1-c^4/2)$, with battery energy and correlations
retained. A static 512-position controller executes the two cycles on a
$32768$-dimensional reducing carrier for all three charged rest energies and
two preparation strengths. A rational Taylor remainder and explicit arithmetic
bounds enclose propagation, reset failure, record overwrite and both flight
phase increments. The second flight has a resolved finite controller bias;
exact preparation-strength independence applies to the isolated collision
identity, not to this complete finite detector. The supplied phase references,
clock packet, terminal measurement work and orthogonal native outcomes are
retained. Sustained replenishment, physical clock units and native twelve-port
execution of this protocol remain separate requirements.

\paragraph{Three-refinement local gauge memory.}
On open full-Haar $SU(2)$ cubes subdivided by $L=1,2,3$, take
$\alpha=32/L$, $\beta=L^3/32$ and a fixed physical unit-square perimeter
of $4L$ links. The Haar probe is supplied, rather than an interacting
vacuum. Fundamental and spin-one Wilson characters have variances $1/4$
and $1/9$ and exact electric frequencies 96 and 256 at every subdivision.
Electric evolution preserves support; a local Duhamel comparison gives
\[
 \delta_L=\min\{2,2\beta N_{\rm inc}|t|\},\qquad
 N_{\rm inc}=5\ (L=1),\quad N_{\rm inc}=8L-4\ (L\ge2).
\]
Normalized real correlation errors are at most $2\delta_L$ and
$3\delta_L$. Both correlations exceed 0.97 at common time $t=1/4096$,
without a global low-energy projector or representation truncation.
The probe is nonstationary, so its normalized correlation is not confined
to $[-1,1]$ by a ground-state argument.

For a Hermitian contraction $A$, the complete instrument
$K_\pm=(I\pm\eta A)/2$, $K_0=\sqrt{(I-\eta^2A^2)/2}$, $\eta=1/2$,
has all nine two-time records. Its weighted moment is exactly
$\mathbb E[s_0s_t]=\eta^2\operatorname{Re}\operatorname{Tr}(\rho A(t)A)$,
including backaction and arbitrary idle references. The Haar derivative
identity bounds the two-readout work by $32/3$ and $512/27$ source units.
Probe preparation and intervention times remain explicit resources.
At fixed time the local error grows as $L^4$; controlled windows of
order $L^{-4}$ supply no fixed-time continuum or Clay theorem
(\path{model/openwave_local_gauge_memory}).

\paragraph{Quantum acquisition of the finite tensor stencil.}
All 60 columns of the native coordinate Wilson source carry ten source
observables: energy, three momenta and six coframe action derivatives.
Three compact declared metric profiles give a full 34-site staggered
stress/curvature stencil. Ten same-state mean records suffice for every
component. Complete three-outcome instruments retain thirty exact
probabilities, the zero events and arbitrary quantum references.
Their nonselective trace-norm disturbance is at most $3/8$, and source
work is at most $3\|H\|/8$ per acquisition. Independent preparations
prevent incompatible sequential readouts from changing the target state.
A two-state control changes $\langle X\rangle$ from 1 to $3/4$ after a
prior $Z$ acquisition, exposing that distinction.

If $C$ is the maximum coefficient norm and each native mean has error
$\epsilon$, every null contraction has error at most $8C\epsilon$.
The sharp inverse $2+\sqrt3$ and measured variation budget yield
\[
 \delta=(2+\sqrt3)(\widehat q+8C\epsilon)
        +\widehat\tau+2C\epsilon,\qquad
 |\Delta\lambda|\le\widehat w+8C\epsilon+8\delta
\]
at unit spacing. The complete sufficient Hoeffding budget is
$3{,}239{,}348{,}495$ independently prepared copies for $\epsilon=10^{-3}$
and joint confidence $0.99$. Every support point and zero halo is retained.
This closes the finite acquisition interface for the stated carrier and
readout assignment; its actual nonzero Einstein residual and physical
metric/stress identification remain explicit
(\path{model/openwave_stress_curvature_instruments}).

\paragraph{Narrow Lambda weight and autonomous quark transport.}
For the actual native weak current applied to a normalizable Lambda packet,
an approximate inverse of the complete interacting Hamiltonian gives
\(\|1_{[E-w,E+w]}(H)u\|\leq\|u-(H-E)x\|+w\|x\|\).
The rational witness includes native proton and pion ground trials at eight
relative separations. Every occupation successor and translation return
contributes to the residual. The conditional weight in the window of
half-width \(1\,\mathrm{MeV}\) is below \(0.026622\), compared with the
moment upper bound \(0.048210\). These are complete finite-resolution
spectral bounds, with the inherited electron and strong calibration.
They imply neither an on-shell density nor a physical decay rate.

The complete source and records are in \path{model/openwave_lambda_narrow_weight}.

Gauge-covariant native hopping transports a fractional-charge endpoint with
its complete reference-supported string. A full-carrier intertwiner reduces
the specified preparation to
\[H|x\rangle=(2+\sigma x)|x\rangle-\kappa(|x-1\rangle+|x+1\rangle).\]
where \(\kappa=1/8\) and \(\sigma=q^2/32+1/12\).
Rational exponential remainders enclose all charge, flux and pointer records
on five- and nine-vertex carriers. The entire fermion and rotor complements
remain in the source. Energy conservation gives the all-time conditional
bound \(\langle x\rangle-1\leq2\kappa/\sigma\), equal to \(6/5\) and
\(24/11\) for native charges \(2/3\) and \(-1/3\).
This autonomous source-unit transport theorem supplies no physical quark
rest-mass or continuum identification.

The source and complete records are in \path{model/openwave_quark_string_transport}.

\paragraph{Stable nuclear attraction and full-input work.}
A neutral harmonic field coupled to a gauge-invariant four-site occupation
projector supplies a signed nuclear interaction while preserving every
sector with fewer than twelve quarks. Completing the full operator square
gives \(H\geq H_{\rm EM}-g^2\sum_b O_b/(2\omega)\), an extensive lower
bound. Exact coherent-state trials provide the corresponding complete-source
variational upper bounds; the oscillator is not truncated. With declared
\(g=1/8\) and \(\omega=1/16\), the helion and alpha comparisons fail.
No held-out mass selects either coefficient. For the original full singlet
input, an independent Casimir trace gives electric preparation energies
\(8955/9152\) and \(2949/1360\) source units for three and four cells,
respectively, including all charge sectors.

The alpha interval is $[3.739060107335,\allowbreak3.739677967934]$ GeV;
helion remains in $[2.807625870510,\allowbreak2.808017057294]$ GeV.
The necessary $|g|\ge1.694976066305$ at fixed oscillator frequency is
a diagnostic obstruction, not a chosen coefficient or sufficient fit.
All oscillator and integer-flux complements remain
(\path{model/openwave_nuclear_field_attraction}).

\paragraph{Complete covariance measurements and normalization.}
For $Q=(q/N)\sum_i Z_i$ in the product $X$ vacuum, two complete coordinate
records $x,y$ give the unbiased statistic
$D=(x-y)(e^x-e^y)/2$ for $C=\operatorname{Cov}(Q,e^Q)$.
With $v=q^2/N$, its second moment is bounded by $3e^{2q}v^2$ and
$C\geq e^{-q}v$. A median of independent block means therefore has a
certified relative-error sample bound independent of $N$. Its measurement
work per pair is $N\omega u$, and its coordinate resolution is a supplied
resource. The complete source gives the weak rest slope
$M(1)=\langle e^Q\rangle$ and the photon slope $M(f_\gamma)$.
Using one clock and length yields
$G_{\rm weak}/(c_\gamma/c)=S C/[2\pi\omega u^2M(1)M(f_\gamma)]$.
For a far photon reference, $M(0)=1$ and this reduces to
$S q\tanh(q/N)/(2\pi\omega u^2)$.
The register count requires independent observational ancestry; the
calibrated count does not predict Newton strength. Finite-source physical
identification and the native-process implementation are not established
by this covariance protocol.

The same pair data also estimate $M(1)$, and a union bound propagates both
relative-error intervals through the complete unit map. No independence
between the two estimators is assumed. Geometry-copy counts for covariance
estimation and local-unbiased source-energy estimation concern distinct
parameters. The former can be independent of $N$ while its detector
resolution and work remain $N$ dependent
(\path{model/openwave_geometry_covariance_protocol}).

\paragraph{Finite-wavelength cosmological transfer.}
For the supplied Einstein--dust--radiation--capacity-scalar model, put
$U=U_m$, $v=U_r-U_m$, $R=\Phi+U$ and $w=y-sU$.
These comoving variables remove the exact adiabatic zero mode. At
$N=\log a=-1$, the declared correction
$\Delta_m=-2Q\Phi/(3\Omega_m)$ enforces the complete Hamiltonian
constraint. An eight-coordinate affine interval equation transports the
$k^2$ response to $N=0$; the exact variation-of-constants equation retains
the complete $k^4$ remainder. The certified metric and comoving-dust
coefficients lie in
\[
 \Phi_1\in[-0.008281806997,-0.005759370982],\qquad
 \Delta_{m,1}\in[-1.455879943233,-1.454618725225].
\]
Full finite-wavelength records are enclosed for $k/H_0=1/100,1/50$,
with overlapping independent 1024- and 2048-step certificates.
The supplied late preparation is distinct from a primordial continuation.
Gravity, density fractions, horizon choice and the common physical unit
remain input premises; no early-Universe abundance or CMB match follows
(\path{model/openwave_capacity_finite_wavelength}).

\paragraph{Native-process admission.}
OpenWave lists OPH as M10 with all 31 criteria untested. Finite theorems,
calibrated extensions and implementation replays are supporting results.
They do not substitute for an in-platform native carrier process generating
the criterion's observable without installing the target field law.
The admission gate preserves this distinction
(\path{model/openwave_m10_review_gate}).
